\documentclass[11pt]{article}
\usepackage[margin=1in]{geometry}
\usepackage{amsmath,amssymb,amsthm}
\usepackage{tikz}
\usetikzlibrary{arrows.meta, positioning, shapes.multipart, fit, backgrounds}
\usepackage{listings}
\usepackage{hyperref}
\usepackage{booktabs}
\usepackage{caption}
\usepackage{graphicx}

\title{Supplementary Note: Balanced-tree derivation, GPU pseudocode, and reproducibility}
\author{STELAR-X supplementary material}
\date{}

\begin{document}
\maketitle

\section*{A. Balanced-tree runtime bound (formal derivation)}

Assume $k$ gene trees, each a perfectly balanced rooted binary tree on $n=2^d$ leaves.
Let $\mathcal{B}$ denote the set of subtree-bipartitions across all gene trees; in the worst case $|\mathcal{B}|=\Theta(nk)$.

\paragraph{Per-pair cost.}
Consider bipartitions $x=(A_1\mid B_1)$ from tree $g_p$ and $y=(A_2\mid B_2)$ from $g_q$.
Let child-subtree sizes be $|A_1|=|B_1|=s$ and $|A_2|=|B_2|=t$.
To compute $M(x,y)$ we need four intersection counts, each computable in $O(\min(s,t))$ time by iterating over the smaller cluster and checking membership via index maps. Hence
\begin{equation}\label{eq:cost}
\mathrm{cost}(M(x,y)) = O(\min(s,t)).
\end{equation}

\paragraph{Aggregate over size classes.}
For size $s=2^i$, the number of bipartitions in one balanced tree with this part-size is $n/(2s)$.
Hence the number of ordered pairs with sizes $(s,t)$ is
\[
\frac{n}{2s}\cdot\frac{n}{2t}.
\]
Their contribution to total cost is
\[
\frac{n}{2s}\cdot\frac{n}{2t}\cdot O(\min(s,t)) = O\!\left(\frac{n^2\min(s,t)}{4st}\right) = O\!\left(\frac{n^2}{4\max(s,t)}\right).
\]
Summing over $s=2^i, t=2^j$ with $0\le i,j\le d-1$ yields
\[
\sum_{i=0}^{d-1}\sum_{j=0}^{d-1} O\!\left(\frac{n^2}{2^{\max(i,j)}}\right) = O(n^2),
\]
because the double sum is bounded by a constant independent of $d$ (numerically $\le 6$).

\paragraph{Over all tree pairs.}
There are $k^2$ ordered tree pairs, yielding total cost $O(n^2 k^2)$ for full precomputation under the balanced-tree assumption.

\section*{B. GPU kernel pseudocode (annotated)}

The following pseudocode outlines the tiled GPU strategy used to precompute $M(x,y)$ contributions and reduce per-$x$ weights. This is intended as a clear, architecture-agnostic description for implementers.

\begin{lstlisting}[basicstyle=\ttfamily\small,caption={Host-side batching and device kernel sketch.}]
// Host: prepare contiguous arrays of tuples for X and Y, prefix scans, and freqY.
for each x_tile in tile_X:
  for each y_tile in tile_Y:
    transfer x_tile and y_tile and related prefix scans to device
    launch kernel compute_tile_contrib(x_tile, y_tile, freqY_tile, out_wx_tile)
    // asynchronously stream next tile to hide PCIe latency
  end
end
// host: aggregate tiled out_wx into global weight map

// Device kernel (each block handles one x entry, threads iterate y entries)
__global__ void compute_tile_contrib(Tuple *tileX, int nx, Tuple *tileY, int ny,
                                    int *freqY, index_maps, prefix_arrays, uint64_t *out_wx) {
  int ix = blockIdx.x;
  if (ix >= nx) return;
  Tuple tx = tileX[ix];
  uint64_t local_sum = 0;
  for (int j = threadIdx.x; j < ny; j += blockDim.x) {
    Tuple ty = tileY[j];
    int f = freqY[j];
    int iA = intersect_count(tx.A, ty.A, index_maps);
    int iB = intersect_count(tx.A, ty.B, index_maps);
    int iC = intersect_count(tx.B, ty.A, index_maps);
    int iD = intersect_count(tx.B, ty.B, index_maps);
    int M = NT(iA,iD) + NT(iB,iC);
    local_sum += uint64_t(M) * uint64_t(f);
  }
  uint64_t block_sum = block_reduce_sum(local_sum);
  if (threadIdx.x == 0) atomicAdd(&out_wx[ix], block_sum);
}
\end{lstlisting}

\paragraph{Implementation notes.}
\begin{itemize}
\item Use 64-bit accumulators for sums to avoid overflow for large $n,k$.
\item Tune tile sizes to fit GPU memory; typical starting tiles: $T_x=1024$, $T_y=4096$ for RTX 4090.
\item Use shared memory to cache tileY when multiple x entries reuse it within a block.
\end{itemize}

\section*{C. Reproducibility and parameter choices}

\begin{table}[h!]\centering
\caption{Recommended parameter defaults for reproducibility.}\label{tab:params}
\begin{tabular}{l l}
\toprule
Parameter & Recommended value \\
\midrule
Field prime $P$ & $2^{61}-1$ (2305843009213693951) \\
Element hash $H(d)$ & $(a\cdot d + b) \bmod P$, salts $a,b$ random \\
Set-hash $\phi_1$ & modular sum in $\mathbb{F}_P$ (associative) \\
Set-hash $\phi_2$ & 64-bit XOR of dispersed $H(d)$ values \\
Tile sizes (GPU) & $T_x=1024$, $T_y=4096$ (adjust per VRAM) \\
CUDA & CUDA 11/12, use thrust/CUB for reductions \\
Fallback CPU mode & Multi-threaded tiling with per-thread local accumulators \\
\bottomrule
\end{tabular}
\end{table}

\bigskip
\noindent\textbf{Hardware used in experiments.} NVIDIA RTX 4090 (24GB VRAM); AMD EPYC 7742 (32 cores); 128 GB RAM.

\bigskip
\noindent\textbf{Software.} CUDA 11/12, C/C++ (GCC 9+), Python 3.8+ for pre- and post-processing scripts.

\end{document}
